\clearpage
\section{Exact padding reflection and contextual congruence}
\label{sec:copying-padding-reflection}
The concrete combined execution now satisfies both representation laws:
recursive-isomorphism invariance and padding invariance. Every finite run
from a padded state reflects to a run of exactly the same length at the
original allowance. Its endpoint is recursively isomorphic to the padding
of the smaller endpoint. Together with the preceding section's forward
transport, this establishes equality of all successful execution costs.
No determinism, refocusing or matching-independence assumption is used.

This is a substantial machine obligation, not the final template theorem.
The abstract cofinal-allowance and full term-context-congruence results can
now be instantiated with this concrete execution without representation-law
premises. The full copying-machine correspondence, matching independence,
refocusing, specimen laws and eleven template families still require proofs.
In particular, this section does not assert equivalence of the two programs
of Proposition~12.2. Their terms and their Figure~16 translations are unchanged.

The dependency chain is: typed truncation and its two round trips;
reflection of stability and copyability; descent of every primitive local
match; descent of contextual matches; exact finite-run reflection; and the
unbounded observation and congruence consequences. The depth invariant is
proved simultaneously with operational descent, rather than assumed.

\Needspace{12\baselineskip}
\subsection{Typed truncation and the representable fragment}

\termdef{depth-truncation}
Write $P_d$ for padding from allowance $d$ to $d+1$, and $L_d$ for the new
auxiliary truncation in the reverse direction. Truncation preserves every
outer vertex, edge, wire label and ordered incidence list. It preserves
primitive labels. At a positive target allowance, it preserves a box's
bound, free-variable and free-atom counts and recursively truncates its body.
At target allowance zero there is no box constructor. To make the auxiliary
function total, an excessive box is replaced by a typed hole with input
$[\mathsf{thunk}(b)]$, output the original free interface, and a fixed
auxiliary name. Such a replacement is used only in the structural proof.
It is neither an execution step nor a clause of term translation.

Totality matters because a context's replaced slot disappears during plugging:
its label need not be representable even if the filled graph is. The proofs
below exclude excessive boxes from actual representable source and target
graphs. No semantic result is claimed for truncating an arbitrary excessive
box without that restriction.

\lean{Copying/DepthLowering}{13}{14}{Lowering.slot}

\lean{Copying/DepthLowering}{17}{25}{Lowering.lower}

\lean{Copying/DepthLowering}{28}{32}{Lowering.label}

The input and output signature equations follow by cases on the target
allowance and on whether the label is primitive or boxed. In the excessive
case, the placeholder has exactly the box's old signature by construction.
Uniform graph truncation is therefore the existing typed relabelling
operation; its graph laws, incidence coverage and finite enumerations are
inherited from that operation. The uniform presentation and recursive
presentation are equal by cases on allowance and labels. The uniform form
makes unchanged outer carriers available definitionally in token proofs.

\lean{Copying/DepthLowering}{34}{36}{Lowering.label_input}

\lean{Copying/DepthLowering}{38}{40}{Lowering.label_output}

\lean{Copying/DepthLowering}{43}{44}{Lowering.graph}

\lean{Copying/DepthLowering}{46}{47}{Lowering.graph_eq_lower}

\termdef{lowering-representability}
A label at allowance $d+1$ \emph{fits} at allowance $d$ when it is primitive,
or when $d>0$ and every label in its boxed body fits one layer lower.
At target allowance zero a box never fits. A graph fits when all its outer
labels do. Thus this is a recursive depth bound on the actual graph, not an
assumption that its syntactic allowance is minimal.

Representability is invariant under recursive isomorphism. At the bottom
layer equivalent labels have the same primitive/box case. At higher layers,
related boxes have identical metadata and a body isomorphism; transport each
body edge using its edge bijection and apply induction. The graph version
uses the outer edge bijection, including its inverse for the forward
universal quantifier. Every padded label fits by induction on allowance:
a padded box has a padded body. Quantifying this label result gives both
the uniform and original graph-padding certificates.

\lean{Copying/DepthLowering}{50}{53}{Lowering.FitsLabel}

\lean{Copying/DepthLowering}{56}{56}{Lowering.Fits}

\lean{Copying/DepthLowering}{59}{86}{Lowering.fitsLabel_equiv}

\lean{Copying/DepthLowering}{88}{95}{Lowering.fits_iso_iff}

\lean{Copying/DepthLowering}{98}{108}{Lowering.fits_padding_label}

\lean{Copying/DepthLowering}{110}{111}{Lowering.fits_graph}

\lean{Copying/DepthLowering}{113}{114}{Lowering.fits_raise}

\Needspace{12\baselineskip}
\subsection{The two round trips and reflection of isomorphism}

Truncation preserves recursive label equivalence even before restricting
to representable graphs. Related primitives remain identical. Related boxes
have the same metadata; at target zero their typed placeholders are identical,
while at a higher target the induction hypothesis transports the body
isomorphism. Consequently graph truncation transports an isomorphism using
exactly its old vertex and edge bijections. Ordered incidence is unchanged.

\lean{Copying/DepthLowering}{117}{147}{Lowering.label_equiv}

\lean{Copying/DepthLowering}{150}{157}{Lowering.graphIso}

Lowering a padded label recovers the original label up to recursive
isomorphism. The primitive case is immediate. In the box case, apply induction
to each body label and use identity outer vertex and edge bijections. The
resulting graph round trip $L_d(P_d(G))\cong G$ also fixes every outer port.
For the converse, $P_d(L_d(H))\cong H$, a representability certificate rules
out the excessive-box case and supplies the induction hypotheses for every
nested body. Its outer bijections are again identities. These proofs establish
that the auxiliary placeholder does not alter representable graphs.

\lean{Copying/DepthLowering}{160}{176}{Lowering.label_roundtrip}

\lean{Copying/DepthLowering}{179}{185}{Lowering.roundtrip}

\lean{Copying/DepthLowering}{188}{206}{Lowering.label_recover}

\lean{Copying/DepthLowering}{209}{216}{Lowering.recover}

\termdef{padding-image}
Representability exactly characterizes the isomorphism-closed image of padding:
\[
 \mathsf{Fits}(H)\quad\Longleftrightarrow\quad
 \exists G\text{ at allowance }d,\ P_d(G)\cong H.
\]
For the forward implication choose $G=L_d(H)$ and use recovery. For the reverse
implication, padded graphs fit and representability transports through the
specified isomorphism. Padding is also faithful on recursive isomorphism:
from $P_d(G)\cong P_d(H)$, lower the isomorphism and compose with the two left
round trips to obtain $G\cong H$. The converse uses the established padding
congruence. No quotient or new graph equality is introduced.

\lean{Copying/DepthLowering}{219}{223}{Lowering.fits_iff_image}

\lean{Copying/DepthLowering}{226}{228}{Lowering.reflectIso}

\lean{Copying/DepthLowering}{230}{236}{Lowering.padding_iso_iff}

\Needspace{12\baselineskip}
\subsection{Assembly, erased slots and search coordinates}

Truncation commutes with actual serial assembly up to an isomorphism
that retains each component tag and the shared interface vertices. Its source
and target equations compose the two incidence maps; cases on an internal,
input or output port establish that its attachment is unchanged. This is the
same explicit finite serial construction used by the original translation.

\lean{Copying/LoweringAssembly}{10}{51}{Lowering.serial_graph}

For simultaneous expansion, truncate each filling and cast its boundary
to the signature of the truncated host label. The signature equalities proved
above justify these casts. Inserting such a cast port in the truncated host
reaches exactly the old vertex. Internal ports retain their component tag;
inputs and outputs retain their source or target index. The expansion
isomorphism has identity outer bijections, and its two incidence equations
follow by mapping this attachment identity over each edge's port list.
Single-slot plugging has the same attachment argument. Retained host edges
use their unchanged host embedding; inserted edges use the attachment identity.

\lean{Copying/LoweringAssembly}{54}{58}{Lowering.expandFilling}

\lean{Copying/LoweringAssembly}{61}{68}{Lowering.expand_insert}

\lean{Copying/LoweringAssembly}{71}{90}{Lowering.expand_graph}

\lean{Copying/LoweringAssembly}{93}{101}{Lowering.plug_insert}

\lean{Copying/LoweringAssembly}{103}{129}{Lowering.plug_graph}

A primitive corolla truncates to the same primitive corolla at the
smaller allowance, by cases on allowance. Equality casts of the external
interface commute with truncation by eliminating those equalities.
Representability is unchanged by either the translation's interface casts
or the graph constructor's reboundary operation: neither changes edge labels.
Serial assembly fits exactly when both components fit, by inspecting its
Boolean-tagged edge carrier. Tensor composition has the same equivalence,
using the general tensor label characterization.

\lean{Copying/LoweringAssembly}{133}{134}{Lowering.primitive}

\lean{Copying/LoweringAssembly}{137}{139}{Lowering.castInterfaces}

\lean{Copying/LoweringAssembly}{142}{143}{Lowering.fits_castInterfaces}

\lean{Copying/LoweringAssembly}{145}{146}{Lowering.fits_reboundary}

\lean{Copying/LoweringAssembly}{149}{155}{Lowering.fits_serial}

\lean{Copying/LoweringAssembly}{158}{160}{Lowering.fits_tensor}

\termdef{lowering-erased-slot}
A plugged graph $C[e:=G]$ fits exactly when $G$ fits and every edge of $C$
other than $e$ has a fitting label. There is deliberately no condition on the
erased label of $e$. The proof splits the replacement's edge carrier into
retained host edges and inserted edges. Likewise, simultaneous expansion fits
exactly when every filling fits; every skeleton edge disappears. This avoids
an unjustified restriction on the existing contextual-match relation.

\lean{Copying/LoweringAssembly}{163}{171}{Lowering.fits_plug}

\lean{Copying/LoweringAssembly}{174}{180}{Lowering.fits_expand}

A box over an allowance-zero body cannot fit one layer lower: apply
the proposed certificate to its unique edge to obtain false. A higher box fits
exactly when its body fits, again by inspecting the unique edge. Truncating a
higher box acts recursively on its body. The lambda-box equation also preserves
the existing bound-first interface convention, using the interface-cast law.
Combining these equations with serial assembly and primitive truncation gives
the isomorphism for complete lambda payloads.

\lean{Copying/LoweringAssembly}{183}{184}{Lowering.box_not_fits}

\lean{Copying/LoweringAssembly}{187}{191}{Lowering.fits_box}

\lean{Copying/LoweringAssembly}{194}{195}{Lowering.box}

\lean{Copying/LoweringAssembly}{198}{201}{Lowering.lambdaBox}

\lean{Copying/LoweringAssembly}{204}{208}{Lowering.lambdaPayload}

On fitting labels, having a specified primitive label is equivalent
before and after truncation. At allowance zero the only problematic case,
a box turning into a hole, is excluded by the certificate; at higher allowances
boxes remain boxes. For any predicate on primitives, the recursive all-labels
predicate is therefore equivalent before and after truncation. The proof
induces through box bodies and uses their fitting certificates. Its graph
version quantifies over unchanged edge carriers. In particular, no focus,
hole or non-language label is introduced into a fitting language graph.

Intrinsic search transport uses identity port and edge bijections, this
primitive-recognition equivalence, and unchanged source and target lists.
External-input status is unchanged. The existing exact search-transport
theorem therefore applies in both directions, with the actual token coordinate.

\lean{Copying/LoweringAssembly}{211}{220}{Lowering.primitive_iff}

\lean{Copying/LoweringAssembly}{223}{236}{Lowering.labelAll_iff}

\lean{Copying/LoweringAssembly}{238}{242}{Lowering.allLabels_iff}

\lean{Copying/LoweringAssembly}{245}{251}{Lowering.search}

\Needspace{12\baselineskip}
\subsection{Reflection of the operational side conditions}

Stable graphs fit at every smaller positive-to-predecessor allowance:
each stable edge is a primitive instance or passive operation, never a box.
Truncation preserves its primitive label and stability classification. The
existing theorem deriving the full stability predicate from these labels
recovers the output and reachability clauses as well. Conversely, for a fitting
graph, stable labels after truncation reflect to the original primitive labels,
so the same argument proves stability of the original graph. Stability has not
been weakened to an informal value test.

\lean{Copying/LoweringPredicates}{11}{14}{Stable.fits}

\lean{Copying/LoweringPredicates}{17}{22}{Stable.lower}

\lean{Copying/LoweringPredicates}{24}{32}{Stable.lower_iff}

\termdef{copyability-padding-reflection}
A lambda payload at allowance one genuinely requires its box layer: serial
representability would imply representability of that box, contradicting the
bottom-layer box obstruction. At a higher allowance, a lambda payload fits
exactly when its body does; the lambda head always fits, and the box/interface
lemmas identify its body's condition.

To lower any copyable graph, induct on its copyability derivation, retaining
an equality that its allowance is a successor. The instance, constant,
application, unary and binary cases truncate to their original primitive
corollas. In the lambda case, target allowance zero is impossible by the
obstruction. Otherwise the lambda-payload isomorphism identifies the result
with a lambda over the truncated body, which is copyable by definition.
No stability premise is imposed on that body. In the renaming case, pull
representability back through the isomorphism, use induction, and transport
copyability through the lowered isomorphism.

\lean{Copying/LoweringPredicates}{38}{43}{Copyable.lambda_not_fits}

\lean{Copying/LoweringPredicates}{46}{54}{Copyable.lambda_fits_iff}

\lean{Copying/LoweringPredicates}{58}{90}{Copyable.lower}

For a fitting graph, copyability of its lowering also implies original
copyability: pad the copyable lowering using the established forward theorem
and transport it along recovery. Thus copyability is equivalent on that
fragment. Specializing to $P_d(G)$ and using the left round trip proves
\[
 \mathsf{Copyable}(P_d(G))\ \Longleftrightarrow\ \mathsf{Copyable}(G).
\]
Induction yields the result for any finite number of layers. Equality transport
in the depth index is eliminated directly; combining it with repeated padding
gives the exact pad-to-allowance equivalence. Extra allowance therefore cannot
make a previously non-copyable payload eligible for the copy rule.

\lean{Copying/LoweringPredicates}{93}{100}{Copyable.lower_iff}

\lean{Copying/LoweringPredicates}{103}{109}{Copyable.raise_iff}

\lean{Copying/LoweringPredicates}{111}{115}{Copyable.pad_iff}

\lean{Copying/LoweringPredicates}{117}{118}{Copyable.castDepth_iff}

\lean{Copying/LoweringPredicates}{120}{122}{Copyable.padTo_iff}

\Needspace{12\baselineskip}
\subsection{Every local match descends}

The recovery isomorphism is first restated using the original recursive
raising function. If a specified padded graph is isomorphic to a higher graph,
lowering that isomorphism and composing with the inverse left round trip
identifies the smaller graph with the lowering of the higher graph.

Every local rule uses an ordered external input as its active port, which
is fixed by any boundary-preserving graph isomorphism. Consequently a known
smaller input-rooted rewrite transports through specified padded source and
target isomorphisms to a rewrite between their lowerings. The target
isomorphism also certifies representability of the higher target, because
all padded graphs fit. This helper supplies both conclusions required by
operational induction.

\lean{Copying/LocalLowering}{13}{15}{Lowering.recover_raise}

\lean{Copying/LocalLowering}{18}{20}{Lowering.fromRaised}

\lean{Copying/LocalLowering}{27}{32}{Combined.lower_input_rule}

\termdef{local-match-lowering}
Suppose an actual primitive local rule takes $(G,v)$ to $(H,w)$ at allowance
$d+1$ and $G$ fits at $d$. Then $H$ fits at $d$ and the actual contextual rewrite
relation at allowance $d$ relates $(L_d(G),v)$ to $(L_d(H),w)$.
The proof covers every constructor, as follows.
\begin{enumerate}
\item Allocation and assignment lower their stable argument. Stable graphs
fit automatically. Their existing redex/reduct padding isomorphisms, followed
by argument recovery, compare the smaller local rule to the given higher one.
Assignment retains exactly the matched atom tree.
\item Dereference and both atom-equality cases have entirely primitive
local graphs. Their original padding equalities directly supply the smaller
rule, with the same matched atom trees and indexed active input.
\item Micro-beta extracts representability of the closure from the source's
serial and tensor decomposition. At target zero its lambda box contradicts
representability. Otherwise the lambda-body certificate and stable-argument
certificate recover both smaller operands. The established before/after
padding and congruence isomorphisms compare the two actual beta rules.
\item Copying extracts the payload certificate from the source's serial
decomposition. Copyability lowers by the new converse theorem; the existing
copy-construction congruence and padding isomorphisms retain the same
contraction-tree presentation on both sides. This is not yet a proof of
independence between different possible tree matches.
\item Unary and binary arithmetic use their established redex-padding and
result-number isomorphisms. Their operands, operation codes and numeric
results are unchanged.
\end{enumerate}
All cases use the preceding input-rooted helper. No local rule is added to
make this proof succeed, and no source or target relation is redefined.

\lean{Copying/LocalLowering}{36}{96}{Combined.local_lowering}

\Needspace{12\baselineskip}
\subsection{Contextual matches and exact finite runs}

Induct on the existing contextual rewrite derivation. A local step is
the preceding theorem. For recursive renaming, pull source representability
back, apply induction, and push target representability and the lowered
rewrite forward. For plugging, the exact representability characterization
provides both the inserted-source certificate and the unchanged surviving-host
certificates. Induction supplies the inserted-target certificate. The new
attachment equations and plugging isomorphisms transport the lower rewrite
at precisely the original inserted token coordinates. For simultaneous
expansion, every source filling fits. Only the selected filling changes;
induction handles it and equality handles the others. The expansion attachment
law then transports the rewrite. This covers arbitrary nested derivations of
the shallow contextual relation, not just canonical local redexes.

\lean{Copying/CombinedLowering}{11}{51}{Combined.rewrite_lowering}

A lowered state consists of the truncated graph and the same mode and
vertex. Its finality condition is literally unchanged: backtracking at an
external input. A search step lowers by the concrete search transport. A
rewrite step lowers by contextual match descent, which also proves that its
target still fits. Induction on finite runs therefore proves both endpoint
representability and a run of exactly the original length between the
lowered endpoints. The induction carries the source certificate through each
step; no global preservation assertion is left as a hypothesis.

\lean{Copying/CombinedLowering}{54}{54}{Combined.lowerState}

\lean{Copying/CombinedLowering}{56}{57}{Combined.lowerState_final}

\lean{Copying/CombinedLowering}{60}{67}{Combined.step_lowering}

\lean{Copying/CombinedLowering}{70}{78}{Combined.run_lowering}

The two graph round trips lift to state isomorphisms because their
outer bijections are identities and thus fix even internal token vertices.
An arbitrary state isomorphism also pads, using the already-proved graph
padding transport while retaining its mode and vertex equations.

\lean{Copying/CombinedLowering}{81}{84}{Combined.lower_pad_state}

\lean{Copying/CombinedLowering}{87}{91}{Combined.pad_lower_state}

\lean{Copying/CombinedLowering}{94}{98}{Combined.padStateIso}

\termdef{execution-padding-reflection}
Let $s$ have allowance $d$, and suppose $P_d(s)\steps{n}t$. Lower the entire
run using the initial padded graph's certificate. This gives a run from
$L_d(P_d(s))$ to $L_d(t)$ and proves that $t$ fits. Transport its starting state
through the left round trip to $s$, obtaining a run $s\steps{n}u$ and a state
isomorphism $L_d(t)\cong u$. Pad that isomorphism, reverse it, and compose with
recovery of $t$ to obtain $P_d(u)\cong t$. Hence
\[
 P_d(s)\steps{n}t\quad\Longrightarrow\quad
 \exists u,\ s\steps{n}u\ \land\ P_d(u)\cong t.
\]
There is no assumption that a rewrite match or a successor is unique.
Finality transports through the endpoint isomorphism. Combining reflection
with forward padding proves equivalence of successful termination at each
exact cost $n$. The result includes zero-step runs.

\lean{Copying/CombinedLowering}{102}{107}{Combined.run_reflect_padding}

\lean{Copying/CombinedLowering}{110}{116}{Combined.terminatesIn_padding_iff}

\Needspace{12\baselineskip}
\subsection{Unbounded executions and observations}

An unbounded run retains its original allowance. The existing bounded-run
extraction theorem therefore turns a run from a padded unbounded state into a
fixed-allowance run. Reflect that run by the preceding theorem and reinsert it
into the unbounded machine. Its endpoint representation relation is obtained
from the reflected endpoint state isomorphism.

For any cost preorder $Q$, every successful higher run now has a successful
lower run at the identical cost, which is related to itself by reflexivity of
$Q$. This proves reverse quantitative refinement. Pairing it with forward
refinement gives padding equivalence for every $Q$, not only the universal
termination preorder. Induction and transitivity give initial-program
equivalence under any finite padding. Specializing to one layer and the
universal preorder discharges the concrete \texttt{RespectsPadding} interface.

\lean{Copying/CombinedPaddingInvariant}{9}{15}{Combined.unbounded_run_reflect_padding}

\lean{Copying/CombinedPaddingInvariant}{18}{22}{Combined.padding_refines_reverse}

\lean{Copying/CombinedPaddingInvariant}{25}{27}{Combined.padding_equiv}

\lean{Copying/CombinedPaddingInvariant}{30}{37}{Combined.initial_pad_equiv}

\lean{Copying/CombinedPaddingInvariant}{40}{41}{Combined.execution_respectsPadding}

\termdef{successful-cost-spectrum}
The successful-cost spectrum of a closed program is the set of natural numbers
$n$ for which its initial state has an $n$-step run to a final state. This set
is equal before and after any finite padding. To prove either inclusion, apply
the corresponding refinement with the equality cost order and eliminate the
resulting equality of costs. Existentially quantifying the cost gives
termination equivalence. Negating that equivalence also preserves
nontermination; no classification of divergent executions is needed.

\lean{Copying/CombinedPaddingInvariant}{44}{53}{Combined.initial_cost_pad_iff}

\lean{Copying/CombinedPaddingInvariant}{56}{62}{Combined.terminates_pad_iff}

For a closing context and an open fragment, first apply initial padding
equivalence to the original closed filling. The established context-padding
isomorphism identifies its padded graph with the filling of the padded context
by the padded fragment. Concrete initial-isomorphism equivalence then gives
exact-cost equivalence of those closed observations. The general observation
reflection theorem can now be instantiated: contextual equivalence at a
larger allowance reflects to the original allowance with both execution
representation premises discharged.

\lean{Copying/CombinedPaddingInvariant}{65}{72}{Combined.close_pad_equiv}

\lean{Copying/CombinedPaddingInvariant}{75}{78}{Combined.equiv_of_pad}

\termdef{concrete-term-congruence}
The cofinal-allowance theorem now applies to the concrete combined execution:
observational equivalence of terms is equivalent to contextual equivalence at
arbitrarily large common adequate allowances. The existing proof includes all
original certified closing contexts by padding them, and now both needed
representation laws are established for the actual machine.
Full congruence now applies to every original annotated term context,
including terms with different minimum translation depths. The proof supplies
concrete isomorphism and padding invariance to the general context theorem.
Its premise is the terms' observational equivalence; the equivalence of
Proposition~12.2 remains to be established.

\lean{Copying/CombinedPaddingInvariant}{81}{85}{Combined.observationalEquiv_iff_cofinal}

\lean{Copying/CombinedPaddingInvariant}{89}{92}{Combined.observationalEquiv_context}

\Needspace{12\baselineskip}
\subsection{Regressions, scope and remaining work}

The bottom-layer regression rejects a lambda payload as representable
without its only box layer. The positive nested-payload regression allows an
arbitrary body and any finite extra allowance, deriving copyability from the
new equivalence. The erased-slot regression allows a nonrepresentable box as
a host slot when a fitting graph replaces it: the host has only one edge, so
there are no surviving host labels to check. The unfilled host still fails
representability. These distinguish actual box depth from irrelevant erased
skeleton labels.

\lean{Copying/LoweringExamples}{10}{11}{Lowering.Examples.lambda_needs_layer}

\lean{Copying/LoweringExamples}{14}{17}{Lowering.Examples.nested_payload_copyable}

\lean{Copying/LoweringExamples}{20}{26}{Lowering.Examples.erased_box_slot_fits}

For every scoped term, padding followed by truncation satisfies the
original translation-language predicate. This uses recursive all-labels
reflection and the existing translation and padding safety proofs.
For each of the two unchanged Proposition~12.2 terms, every certified closing
context at its original minimum allowance and every finite padding preserve
all successful costs. The environments are arbitrary variable and atom counts;
there is no empty-environment restriction. The two assertions compare each
program with its own padded representation, not with the other program.
A final regression preserves nontermination of every closed contextual filling
under padding. Its universal quantification includes admissible contexts
with divergent callbacks or repeated closure use, but does not independently
prove the behaviour of a particular such callback.

\lean{Copying/LoweringExamples}{29}{33}{Lowering.Examples.translated_roundtrip_language}

\lean{Copying/LoweringExamples}{37}{42}{Lowering.Examples.privateCell_observation_padding}

\lean{Copying/LoweringExamples}{45}{50}{Lowering.Examples.constantResult_observation_padding}

\lean{Copying/LoweringExamples}{55}{59}{Lowering.Examples.nontermination_padding}

All 89 new declarations are included in the copying axiom audit. The
full build checks 362 jobs; all three audits pass with 6 abstract, 477 graph
and 1001 copying declarations, using only \texttt{propext},
\texttt{Classical.choice} and \texttt{Quot.sound}. There are no project axioms
or unfinished proofs in this milestone.

Appendix~B.1 no longer lists padding reflection or its observation and
congruence consequences as unfinished. Matching independence, deterministic
behaviour up to representation, refocusing, focused-graph correspondence,
specimen locality and template robustness remain. The completed invariance
law concerns the concrete combined relation as defined; correspondence with
the paper must still be maintained if that relation needs correction.

The updated estimate for Proposition~12.2 is about 78\% of the mathematical
effort, with a broad plausible range of 60--88\% (26 October 2026). The change
from 75\% reflects completion of a substantive semantic interface and its
congruence consequences. This is a judgement about remaining proof difficulty,
not a count of declarations or completed to-do bullets. The eleven-template
argument and final Figure~22 chain retain substantial uncertainty.

